% ===================================================================
%  TABELLA N-RO 9 — Le montania. — Le citate de aquas. — Le foreste.
%  — Le chassa.
%  Traduction 2026 d'apres texto/io, controlee sur texto/fr et
%  texto/es. Consigne : voir texto/ia/10-tabella-01.tex.
% ===================================================================

%%K t09-apar-1 apar
\VUsaut{4.0mm}
\VUtitre{15.0pt}{60}{TABELLA N\textsuperscript{o} 9}
\VUsaut{3.0mm}

%%K t09-tit-1 sub
\VUcentre{12.0pt}{0}{\textit{Prime scena.}}
\VUsaut{3.0mm}

%%K t09-tit-2 sub
\VUpk{11.4pt}{40}{Le montania. --- Le citate de aquas.}
\VUsaut{3.0mm}

%%K t09-c1-01-1 p
1. --- De octo dies io es in Prats. Io non pote exprimer tote le admiration que
io sentiva quando io esseva ante le meraviliose \VUgras{catena de montanias}
\textsuperscript{(1)} del Pyreneos del qual le \VUgras{cresta}
\textsuperscript{(2)} appare sur le blau celo como un gigantesc festone.

%%K t09-c1-02-1 p
2. --- Io non imaginava, que le \VUgras{summitates} pyrenee
\textsuperscript{(3)} attinge un \VUgras{altitude} tanto grande, e io restava
stupefacte ante le belle \VUgras{glaciares} \textsuperscript{(5)} e le
\VUgras{picos} \textsuperscript{(4)} nivose sur le quales rola tanto terribile
\VUgras{avalanchas}.

%%K t09-tit-3 sub
\VUsaut{3.0mm}
\VUpk{10.2pt}{40}{Paisage de montania.}
\VUsaut{3.0mm}

%%K t09-c1-03-1 p
3. --- Ja le die post mi arrivata, io iva al vetule \VUgras{vulcano}
\textsuperscript{(6)} extincte que es presso Prats. Sur le bordos aride e nudate
del \VUgras{crater}, on vide ancora ben le tracias del passage del \VUgras{lava}
que flueva, plure seculos ante ora, sur le \VUgras{flanco del montania}
\textsuperscript{(7)}. Io succedeva trovar, sur un del \VUgras{declivos}, alcun
fragmentos de ille lava.

%%K t09-c1-04-1 p
4. --- Prats es situate sur le frontiera, e le \VUgras{fortalessa}
\textsuperscript{(8)} que defende le \VUgras{gorga} \textsuperscript{(27)} que
separa Francia de Espania rememora completemente ille vetule castellos del ripa
del Rheno qui pare guattar un \VUgras{preda} del summitate de lor roca.

%%K t09-c1-05-1 p
5. --- Un enorme \VUgras{aquila} \textsuperscript{(9)} qui ha su nido non
lontano del \VUgras{forte} \textsuperscript{(8)} sovente attrahe mi attention.
Iste \VUgras{ave rapace}, del qual le \VUgras{becco} es forte, sovente apporta a
su pullos un agno o un capreto que illo tene per su \VUgras{ungues}. Tanto es le
grandor de su \VUgras{alas} que illo pote longemente planar con su pesante
carga.

%%K t09-tit-4 sub
\VUsaut{3.0mm}
\VUpk{10.2pt}{40}{Le excursionista.}
\VUsaut{3.0mm}

%%K t09-c1-06-1 p
6. --- Illac le plus agradabile promenada es le excursion al cascada
\textsuperscript{(10)} de « Cau ». Le \VUgras{cadita de aqua} cade in turbines e
postea dispare como un belle \VUgras{rivetto} in un \VUgras{bosco de abietes}
\textsuperscript{(11)} situate al west del citate. Nos ascendeva tra iste bosco
usque le \VUgras{observatorio meteorologic} \textsuperscript{(12)} e del
summitate del \VUgras{belvedere}, nos percipeva tote le \VUgras{panorama} del
region. Le \VUgras{paisage} es meraviliose, e le \VUgras{natura} pare prodigar
al pais tote le tresores del quales illo dispone.

%%K t09-c1-07-1 p
7. --- Pro descender, nos usava le \VUgras{funicular} \textsuperscript{(13)} e
nos haltava in un parve \VUgras{station} presso le \VUgras{ruinas}
\textsuperscript{(14)} de un vetule \VUgras{castello feudal} olim habitate per
le seniores de Prats.

%%K t09-c1-08-1 p
8. --- Povre castello! Que pensa illo probabilemente vidente in basso le
coquette \VUgras{citate de aquas} \textsuperscript{(15)} que ora es un
\VUgras{station thermal} de moda?

%%K t09-c1-08-2 p
Le \VUgras{clima} es tanto agradabile, le \VUgras{aqua mineral} tanto efficace,
que le \VUgras{cura} usate in le \VUgras{sanatorio} \textsuperscript{(17)} es un
ver placer.

%%K t09-tit-5 sub
\VUsaut{3.0mm}
\VUpk{10.2pt}{40}{Agradabile citate thermal.}
\VUsaut{3.0mm}

%%K t09-c1-09-1 p
9. --- Del resto, on faceva toto pro render agradabile le sojorno. Un grande
\VUgras{viaducto} \textsuperscript{(16)} esseva construite pro permitter un
ferrovia; le \VUgras{parco} \textsuperscript{(18)} del \VUgras{establimento
thermal}, in le qual on da \VUgras{concertos}, es arrangiate in un maniera
incantatori. Plure gratiose « \VUgras{villas} » \textsuperscript{(19)} esseva
construite circa le casino \textsuperscript{(21)}, e un \VUgras{chaussée}
\textsuperscript{(22)} multo ben mantenite grandemente facilita le
communicationes.

%%K t09-c1-10-1 p
10. --- Un simple \VUgras{taluto} \textsuperscript{(23)} separa le \VUgras{via}
\textsuperscript{(20)} del \VUgras{valle} \textsuperscript{(25)} proprie dicte, e
in le fundo del qual jace un gratiose parve \VUgras{laco}
\textsuperscript{(26)} que alimenta un \VUgras{rivetto} \textsuperscript{(24)}
limpide.

%%K t09-c1-11-1 p
11. --- Al altere latere del valle, un \VUgras{mina de ferro}
\textsuperscript{(28)} es exploitate. Plure vices io iva vider le \VUgras{mineros}
descendente in le \VUgras{puteo}, e mesmo io entrava le \VUgras{galeria} in le
qual illes passa lor die, extrahente le \VUgras{mineral} que contine le
\VUgras{metallo}.

%%K t09-c1-12-1 p
12. --- Un importante \VUgras{fonderia} esseva construite presso le mina pro
utilisar immediatemente le ricchessas que on extrahe de illo. Le \VUgras{alte
furno} \textsuperscript{(29)}, calefacite per le \VUgras{carbon de terra} que
abunda hic, permitte obtener rapidemente e a bon mercato le \VUgras{ferro
fundite}.

%%K t09-c1-13-1 p
13. --- Non lontano del mina on excava un \VUgras{petreria}
\textsuperscript{(30)}. Ni \VUgras{granito}, ni \VUgras{porphyro}, ni mesmo
\VUgras{gres} le vetule \VUgras{petrero} talia con su \VUgras{picco}, ma simple
\VUgras{petras de talia} pro le construction de casas.

%%K t09-c1-14-1 p
14. --- Un del plus belle ornamentos de iste pais es le \VUgras{abietes}
\textsuperscript{(31)}. Ubique in le fundo de un \VUgras{ravina}
\textsuperscript{(32)}, sur le ripa de un \VUgras{torrente}
\textsuperscript{(33)}, de un \VUgras{abysmo} \textsuperscript{(34)} o de un
precipitio, lor tenue acos nuancia in verde.

%%K t09-tit-6 sub
\VUsaut{3.0mm}
\VUpk{10.2pt}{40}{Marcha a pede.}
\VUsaut{3.0mm}

%%K t09-c1-15-1 p
15. --- Io profita mi vacantias pro promenar \VUgras{longemente}. Le
\VUgras{vias} \textsuperscript{(35)} qui serpenta sur le montania permitte
percurrer, sin reguardar le distantia, plure \VUgras{kilometros} e sovente plure
\VUgras{myriametros}.

%%K t09-c1-15-2 p
\VUgras{Petras de limite} \textsuperscript{(36)} e \VUgras{postos indicator}
\textsuperscript{(37)} marca le vias sur le quales sovente on incontra un juvene
\VUgras{montaniero} \textsuperscript{(38)} con un \VUgras{bisacco}
\textsuperscript{(40)} a su flanco, portante un \VUgras{marmotta}
\textsuperscript{(39)}, o alcun \VUgras{gitano} \textsuperscript{(41)} del qual
le \VUgras{toque de pelle} \textsuperscript{(42)} contrasta estranimente con le
vetule \VUgras{capuce} \textsuperscript{(43)} lacerate. Ille conduce per un
\VUgras{catena un urso} \textsuperscript{(44)} dressate.

%%K t09-tit-7 sub
\VUpk{10.2pt}{40}{Ascension de montanias.}
\VUsaut{3.0mm}

%%K t09-c1-16-1 p
16. --- Quasi cata die io va con un \VUgras{guida} \textsuperscript{(48)}
practicar un \VUgras{ascension}, e io es multo fer quando, con un \VUgras{sacco}
\textsuperscript{(47)} al dorso, e le « \VUgras{alpenstock} » in le mano, como
un ver \VUgras{turista} \textsuperscript{(46)} io assalta le \VUgras{rocas}
\textsuperscript{(45)}.

%%K t09-c1-17-1 p
17. --- Del summitate de un montania, le habitantes del planitia pare non plus
grosse que formicas; un \VUgras{grege de oves} \textsuperscript{(50)} pasturante
in un \VUgras{pastura} \textsuperscript{(49)} pare un ludo pro infantes. Un
\VUgras{pino} \textsuperscript{(51)} o anque un \VUgras{parve casa de ligno}
\textsuperscript{(52)} appare a pena como un macula sur le verdura.

%%K t09-c1-18-1 p
18. --- Durante iste excursiones, nos sovente incontra sur le \VUgras{sentiero}
\textsuperscript{(53)} un \VUgras{chassator de chamois} \textsuperscript{(54)}
portante sur le spatula le bestia que ille justo occideva.

%%K t09-c1-19-1 p
19. --- Io poneva in mi herbario un multo remarcabile branca de \VUgras{filice}
\textsuperscript{(55)} e un belle branchetta rosate de \VUgras{erica}
\textsuperscript{(57)} que io colligeva al entrata de un parve \VUgras{grotta}
\textsuperscript{(56)} durante mi ultime promenada.

%%K t09-tit-8 sub
\VUsaut{3.0mm}
\VUcentre{12.0pt}{0}{\textit{Secunde scena.}}
\VUsaut{3.0mm}

%%K t09-tit-9 sub
\VUpk{11.4pt}{40}{Le foreste. --- Le chassa.}
\VUsaut{3.0mm}

%%K t09-c2-01-1 p
1. --- Heri io passava un die deliciose in le foreste. Le diligentia que regna
inter le parve animales e le \VUgras{insectos} \textsuperscript{(5)} per le
quales illo es habitate multo me interessava.

%%K t09-c2-01-2 p
Le \VUgras{aranea} \textsuperscript{(1)} qui texe su \VUgras{tela}
\textsuperscript{(2)} inter duo brancas pro capturar le \VUgras{musca}
\textsuperscript{(3)} passante, le \VUgras{caracol} \textsuperscript{(4)}
portante penosemente su concha, le lucano explorante pro su preda, le diligente
formica multo zelose presso le \VUgras{formichiero} \textsuperscript{(6)}, tote
iste parvos iva, veniva, laborava con un extraordinari diligentia.

%%K t09-c2-02-1 p
2. --- Un \VUgras{serpente} \textsuperscript{(7)} que io a prime vista opinava
un \VUgras{vipera}, ma que non esseva altere cosa que un simple
\VUgras{colubro}, se habeva tapite sub un tronco de arbore, e de tempore a
tempore faceva exir su tenue parve capite que sempre pareva preste pro morder.
Ante me un grosse \VUgras{limace} \textsuperscript{(8)} reptava pesantemente post
haber devorate un del sapide \VUgras{fragas} \textsuperscript{(11)} qui cresce
abundantemente illac.

%%K t09-c2-03-1 p
3. --- Le solo esseva tapetate de \VUgras{flores de bosco}; \VUgras{primulas}
\textsuperscript{(9)}, \VUgras{pervincas} \textsuperscript{(10)}, e multe altere
plantas que io non cognosceva, floreva inter \VUgras{tufos de herba}
\textsuperscript{(12)}.

%%K t09-c2-04-1 p
4. --- In iste region, le \VUgras{trufa} es quasi tanto commun como le
\VUgras{fungo} \textsuperscript{(14)} e un \VUgras{porchetto}
\textsuperscript{(13)} appertinente a un forestero, gulosemente explorava si
illo succederea trovar alcunes de illos.

%%K t09-tit-10 sub
\VUsaut{3.0mm}
\VUpk{10.2pt}{40}{Chassa furtive.}
\VUsaut{3.0mm}

%%K t09-c2-05-1 p
5. --- Io assisteva al \VUgras{fin} de un scena que multo me amusava. Con le
adjuta del olfacto de su \VUgras{can basset} \textsuperscript{(15)}, le
\VUgras{forestero} \textsuperscript{(16)} justo surprendeva un \VUgras{chassator
furtive} \textsuperscript{(22)}, al momento quando iste se preparava a portar
via le \VUgras{gibier} \textsuperscript{(17)} que ille habeva occidite.
Preferente abandonar su preda que lassar se arrestar, le chassator furtive habeva
fugite, e un \VUgras{lepore} \textsuperscript{(18)}, un \VUgras{cerva}
\textsuperscript{(19)}, un \VUgras{capreolo} \textsuperscript{(20)} e un
\VUgras{apro} \textsuperscript{(21)} jaceva sur le herba allegrante le forestero.

%%K t09-c2-06-1 p
6. --- Le diversitate del arbores que contine le foreste es stupefaciente. Le
\VUgras{fagos} \textsuperscript{(23)} e le \VUgras{betulas}
\textsuperscript{(26)}, le \VUgras{pinos}, qui produce le \VUgras{resina}
\textsuperscript{(27)}, le \VUgras{quercos} \textsuperscript{(29)} e anque multe
alteres intermisce lor brancage.

%%K t09-c2-06-2 p
Le \VUgras{hedera} \textsuperscript{(24)}, que adhere al tronco del alte
arbores, colora le \VUgras{silva} \textsuperscript{(25)} de un verdura
brillante, que se uni agradabilemente al nuance plus clar e pallidemente verde
del \VUgras{musco} \textsuperscript{(31)} in le qual le \VUgras{pico verde}
\textsuperscript{(30)} cerca insectos.

%%K t09-tit-11 sub
\VUsaut{3.0mm}
\VUpk{10.2pt}{40}{Paisage de foreste.}
\VUsaut{3.0mm}

%%K t09-c2-07-1 p
7. --- Le vetule quercos me incantava. Lor grosse \VUgras{radices}
\textsuperscript{(32)} torte, medio exite del solo, les da un splendide aspecto
de fortia e vigor, e io me demanda como le \VUgras{proprietario}
\textsuperscript{(35)} pote resignar se a facer los abatter e livrar los al
\VUgras{serra} \textsuperscript{(34)} del \VUgras{buchero}
\textsuperscript{(33)}. Le povre \VUgras{truncos} \textsuperscript{(36)}, private
de lor \VUgras{cortice} \textsuperscript{(37)} o findite per le \VUgras{hacha}
\textsuperscript{(38)} del \VUgras{obrero a jornata} \textsuperscript{(39)} me
attrista.

%%K t09-c2-08-1 p
8. --- Presso, un vetule \VUgras{carbonero} \textsuperscript{(41)} preparava con
su \VUgras{pala un amasso} \textsuperscript{(42)} de carbon, e un ancian
portava un \VUgras{fasce} \textsuperscript{(40)} de \VUgras{ligno morte} de
brancas del arbores abattite.

%%K t09-tit-12 sub
\VUsaut{3.0mm}
\VUpk{10.2pt}{40}{Chassa.}
\VUsaut{3.0mm}

%%K t09-c2-09-1 p
9. --- Io voleva ir pro reposar durante alcun momentos in le \VUgras{casa del
forestero} \textsuperscript{(46)} ma, quando io arrivava presso le \VUgras{parve
laco} \textsuperscript{(43)}, sur le qual nata un belle \VUgras{cygno}
\textsuperscript{(45)}, io remarcava, que un recente \VUgras{inundation}
\textsuperscript{(44)} habeva cambiate le via in un ver \VUgras{palude}. Io
debeva deviar me, e, passante presso le \VUgras{cedro} \textsuperscript{(49)},
le \VUgras{poplos} \textsuperscript{(48)}, e le \VUgras{ulmo}
\textsuperscript{(47)} qui es al bordo del foreste, io videva emerger de un
\VUgras{copse} \textsuperscript{(54)} un belle \VUgras{cervo}
\textsuperscript{(50)}, persequite per un obstinate \VUgras{muta de canes}
\textsuperscript{(51)} que excitava anque le \VUgras{corno} \textsuperscript{(53)}
del \VUgras{chassatores a cavallo} \textsuperscript{(52)}. Le povre bestia
esseva toto exhauste. Illo veniva morente cader a alcun passos presso me.

%%K t09-c2-10-1 p
10. --- Le filio del forestero, que le ruito del \VUgras{chassa a curso} habeva
attrahite, exiva del casa e nos duo retornava in le foreste. Ille habeva vidite
un \VUgras{pica} \textsuperscript{(56)} sur un branca de un vetule querco. Ille
opinava, que su nido probabilemente es celate in le \VUgras{foliage}
\textsuperscript{(58)} e ille voleva capturar le nido.

%%K t09-c2-10-2 p
Agile como un \VUgras{skiuro} \textsuperscript{(61)}, ille ascendeva de
\VUgras{branca} \textsuperscript{(55)} a branca, ma ille trovava nihil e
succedeva solmente facer volar un vetule \VUgras{cornice} \textsuperscript{(60)}
posate sur un \VUgras{grande branca} \textsuperscript{(57)}.
\VUsaut{4.0mm}
\VUfilet{22mm}
